\chapter{Experimental Design}
\label{ch:experiments}

\section{Completed technical verification}

Completed verification comprises source/schema audit, file and database consistency checks, automated tests against an isolated database, representative API probes against a disposable database copy, a clean frontend dependency/build run, browser inspection of eight main surfaces, and LaTeX build/render inspection. These procedures answer whether the snapshot is internally reproducible and demonstrable; they do not answer whether users receive relevant or faithful advice.

\section{Implemented retrieval metrics}

For query $q$, with retrieved paper list $L_q$ and gold set $G_q$, the code reports a binary hit indicator called Recall@$k$:
\begin{equation}
\operatorname{Hit@}k(q)=\mathbb{I}\!\left[G_q\cap L_q[1{:}k]\neq\varnothing\right].
\label{eq:hitk}
\end{equation}
The implementation averages this indicator over queries. It is therefore a hit rate, not set recall $|G_q\cap L_q[1{:}k]|/|G_q|$. Mean reciprocal rank is
\begin{equation}
\operatorname{MRR}=\frac{1}{|Q|}\sum_{q\in Q}\frac{1}{\operatorname{rank}_q},
\label{eq:mrr}
\end{equation}
with zero when no gold paper is retrieved. The sample retrieval and QA files contain empty gold labels and intentionally fail validation, so neither metric has a current result.

\section{Proposed retrieval experiment}

A versioned test set should stratify questions by exact fact, method, result, limitation/future work, topic, author, comparison, and broad intent. Keyword, hashing-vector, and hybrid modes should use the same corpus, cut-offs, and paper filtering. Primary outcomes should include set Recall@3/5, hit rate, MRR, nDCG where graded judgments are available, and per-intent error analysis. Query-expansion and reranking terms should be ablated. The active vector index must first be rebuilt for all 756 chunks and status consistency verified.

\section{Proposed QA evaluation}

Each answer should be decomposed into checkable claims. Reviewers should label whether each claim is entailed by its cited snippet, whether the citation points to the right paper/page, whether important answer points are covered, and whether unsupported general knowledge appears. Automatic RAG measures may supplement but not replace these labels \cite{es2024ragas,saadfalcon2024ares,min2023factscore}. Provider, model hash/tag, prompt, temperature, retrieval list, and fallback state must be stored.

\section{Proposed recommendation evaluation}

Students and supervisors should independently rate paper relevance, gap fidelity, novelty, feasibility for the stated time/skills/data, usefulness of the MVP, risk calibration, and appropriateness of the evaluation plan. Reviewers must separately judge source claims and system suggestions. A comparison baseline could present retrieved papers without generated extension advice. \authorinput{approve participant groups, sample size, rating scales, compensation, consent, and analysis}.

\section{Proposed topic/author and summary evaluation}

Topic evaluation should measure label precision/coverage and reviewer agreement against the fixed taxonomy, including multi-label cases and an ``other'' outcome. Author evaluation should first resolve identities, then assess expertise links. Public and technical summaries should be reviewed for correctness, relevance, coherence, and readability, drawing on human-centred evaluation lessons from SummEval \cite{fabbri2021summeval}.

\section{Evaluation workflow and current state}

\begin{figure}[htbp]
\centering
\begingroup\hfuzz=100pt\resizebox{\textwidth}{!}{\begingroup\hfuzz=100pt
\begin{tikzpicture}[node distance=8mm and 10mm]
  \node[flowbox] (questions) {Manually reviewed questions\\and gold paper IDs};
  \node[flowbox, right=of questions] (retrieval) {Run keyword, hashing,\\and hybrid retrieval};
  \node[flowbox, right=of retrieval] (metrics) {Hit-based Recall@3/5\\and MRR};
  \node[flowbox, right=of metrics, fill=PaleGold, draw=Gold!80!black] (missing) {Current state:\\gold labels absent; not run};

  \node[flowbox, below=15mm of questions] (qa) {Reviewed QA cases +\\required answer points};
  \node[flowbox, right=of qa] (outputs) {Answers, recommendations,\\summaries and citations};
  \node[flowbox, right=of outputs] (automatic) {Citation presence, warning\\and structural status counts};
  \node[flowbox, right=of automatic, fill=PaleGold, draw=Gold!80!black] (placeholder) {Current state:\\one placeholder case per\\extension/artifact evaluation};

  \node[flowbox, below=15mm of qa] (rubric) {Human rubric:\\faithfulness, usefulness,\\citation correctness, feasibility};
  \node[flowbox, right=of rubric] (review) {Independent reviewers +\\adjudication and\\reliability};
  \node[flowbox, right=of review] (analysis) {Error taxonomy, intervals,\\comparisons and conclusions};
  \node[flowbox, right=of analysis, fill=PaleGold, draw=Gold!80!black] (notdone) {Current state:\\templates only; zero reviews};

  \draw[arrow] (questions) -- (retrieval);
  \draw[arrow] (retrieval) -- (metrics);
  \draw[arrow] (metrics) -- (missing);
  \draw[arrow] (qa) -- (outputs);
  \draw[arrow] (outputs) -- (automatic);
  \draw[arrow] (automatic) -- (placeholder);
  \draw[dashedarrow] (rubric) -- (review);
  \draw[dashedarrow] (review) -- (analysis);
  \draw[dashedarrow] (analysis) -- (notdone);
\end{tikzpicture}
\endgroup
}\endgroup
\caption{Implemented and proposed evaluation workflow. Gold retrieval/QA labels and human reviews are absent in the audited snapshot.}
\label{fig:evaluation-workflow}
\end{figure}

\begin{figure}[htbp]
\centering
\projectScreenshot{figures/screenshots/evaluation.png}{0.96\textwidth}
\caption{Evaluation dashboard correctly displaying not-run and placeholder states rather than inferred performance.}
\label{fig:evaluation-shot}
\end{figure}
